\begin{helper}
Let \(\eta,\varepsilon_1\) be real numbers with
\[
0<\eta,\qquad 0<\varepsilon_1,\qquad \eta<\frac1{16}.
\]
Then there is \(n_0\) such that, for every \(n>n_0\), if
\[
L=\log n,\qquad k_R=(1+\eta)\sqrt{nL},
\]
then
\[
n^{4\eta}\le\frac{\varepsilon_1}{2}k_R
\qquad\text{and}\qquad
2L^2\le\frac{\varepsilon_1}{8}k_R.
\]
\end{helper}
\begin{proof}
Put
\[
c_1=\frac{\varepsilon_1}{2}(1+\eta),\qquad
c_2=\frac{\varepsilon_1}{8}(1+\eta).
\]
Both constants are positive because \(\varepsilon_1>0\) and
\(1+\eta>0\).  Since \(\eta<1/16\), we have \(4\eta<1/4\).

First choose the threshold so large that \(n\ge1\), \(L=\log n\ge1\), and
\[
1\le c_1 n^{1/4}.
\]
These are eventual conditions: \(n\to\infty\), \(\log n\to\infty\), and
\(c_1n^{1/4}\to\infty\).  For such \(n\), monotonicity of real powers on
\([1,\infty)\) gives
\[
n^{4\eta}\le n^{1/4}.
\]
Multiplying \(1\le c_1n^{1/4}\) by \(n^{1/4}\ge0\) gives
\[
n^{1/4}\le c_1n^{1/2}=c_1\sqrt n.
\]
Since \(L\ge1\), we also have \(n\le nL\), hence
\(\sqrt n\le\sqrt{nL}\).  Combining these inequalities gives
\[
n^{4\eta}\le c_1\sqrt{nL}
=\frac{\varepsilon_1}{2}(1+\eta)\sqrt{nL}
=\frac{\varepsilon_1}{2}k_R.
\]

For the second inequality choose
\[
\delta=\sqrt{c_2/4}>0.
\]
Because \(\log n=o(n^{1/4})\), enlarge the threshold so that, for every
\(n>n_0\),
\[
L=\log n\le\delta n^{1/4}.
\]
Together with \(L\ge0\), squaring gives
\[
L^2\le \delta^2 n^{1/2}=\delta^2\sqrt n.
\]
Since \(\delta^2=c_2/4\), we have \(2\delta^2\le c_2\), and therefore
\[
2L^2\le c_2\sqrt n.
\]
As above, \(L\ge1\) implies \(\sqrt n\le\sqrt{nL}\).  Thus
\[
2L^2\le c_2\sqrt{nL}
=\frac{\varepsilon_1}{8}(1+\eta)\sqrt{nL}
=\frac{\varepsilon_1}{8}k_R.
\]
Taking the maximum of the finitely many thresholds used above proves the
claim.
\end{proof}
